\section{Additional Experiments}
\subsection{Experiment A: Homework-Compliant DETR vs Deformable DETR}
\textbf{Hypothesis.}
With the same rule (only pretrained ResNet-50 backbone, no pretrained transformer), Deformable DETR should converge better than DETR.

\textbf{What may (or may not) work.}
Architecture upgrade should improve feature aggregation and box refinement even without transformer pretraining.

\textbf{Result.}
Comparing successful runs R1 and R2:
\begin{enumerate}
    \item DETR baseline (R1): latest AP=0.00003 (best AP=0.00039), and convergence is very slow.
    \item Deformable DETR (R2): AP=0.2993.
\end{enumerate}

\textbf{Implication.}
Under the assignment constraint, moving from DETR to Deformable DETR is the main source of improvement in my experiments.

\subsection{Experiment B: Learning Progression in a Successful Deformable Run}
\textbf{Hypothesis.}
Without transformer pretraining, the model should still improve steadily over epochs if the training pipeline is correct.

\textbf{What may (or may not) work.}
Early epochs may focus on coarse localization and classification, while later epochs improve precision.

\textbf{Result.}
In R2 (\texttt{kqscm42c}):
\begin{enumerate}
    \item Epoch 1: AP=0.0329.
    \item Epoch 10: AP=0.2833.
    \item Epoch 13: AP=0.2993.
\end{enumerate}

\textbf{Implication.}
The Deformable DETR pipeline converges reliably in this successful setup.

\subsection{Experiment C: Effect on AP50 and AP}
\textbf{Hypothesis.}
If localization quality improves together with classification confidence, AP50 and AP should rise together.

\textbf{What may (or may not) work.}
AP50 may improve faster than AP when coarse localization is learned earlier than precise localization.

\textbf{Result.}
In R2, AP50/AP progressed from 0.0888/0.0329 (epoch 1) to 0.5785/0.2993 (epoch 13).

\textbf{Implication.}
Both metrics improved substantially, indicating better detection quality rather than only threshold-sensitive gains.
